\chapter{Parte Central}
Name: Identificación dinamica de tópicos
[Placeholder]

La parte central del trabajo refiere a lo que es producción propia o aporte del
proyecto de grado, incluyendo las decisiones tomadas. Por ejemplo, puede incluir
los requerimientos, el análisis y el diseño de la solución. Si el proyecto tiene una
implementación, debe describirse en términos de decisiones tomadas en ese sentido.
Los detalles de programación se dejan para los anexos.

Se pueden incluir figuras y tablas en el documento, las mismas deben estar referenciadas en el texto. Por ejemplo, la Figura muestra los logos de Facultad de Ingeniería y de la Universidad de la República, mientras que la Tabla tiene números aleatorios.

[Que poner]

\section{Serie de pasos}

\subsection{Obtención de mensajes}

\subsection{Preprocesamiento de mensajes}

\subsection{Procesamiento de mensajes}

\subsection{Evaluación}

\subsection{Visualizacion de Topicos}


\section{Aplicación}

\subsection{Orquestrador}

\subsection{Usuario}
